\section{Conclusion}

The paper answers the timetable problem through a connected operational chain. First, passenger-flow evidence is used to identify the overlap zone and motivate a large-small-route service design. Second, a combined cost-service criterion is used to determine the selected operation plan. Third, the timetable recurrence model converts that plan into a machine-readable schedule which is then checked through capacity, tracking, and dwell audits.

The core conclusion is that the service pattern and the timetable must be interpreted together. A preferred service pattern is not yet a finished answer unless it can be turned into a feasible timetable, and a feasible timetable is not yet a convincing recommendation unless it can be traced back to the passenger-flow burden that motivated it.

In summary, the proposed route yields not only a selected operation strategy but also an auditable timetable and a scenario-aware operating recommendation. Within the tested perturbation range, the final solution can be regarded as sufficiently feasible and interpretable for research review, while later work should still extend the route with richer OD behavior, vehicle circulation, and more realistic operating constraints.
